\section{Polylogarithmic counterterms for higher Stieltjes collisions}
\label{gcx:section}

The geometric collision theorem at Stieltjes index zero in the incoming
\emph{Higher Reflection Identities and Resonant Zeta Calculus}\cite{repo:reflection} leaves two explicit gaps: logarithmic
Stieltjes singularities, and weighted collision formulas.  The theorem below
covers both, together with arbitrary ordinary argument derivatives.  The
counterterm is finite and explicitly reducible to iterated polylogarithms.
The proof also explains why no comparison between different shrinking gaps
is needed.  It does not evaluate the global merged cubic constant, and it
does not assert associativity of a different, successively chosen
subtraction prescription.

\subsection{Local singular polynomials and the finite-part convention}

Normalize generalized Stieltjes functions by
\[
 \zeta(1+u,a)=\frac1u+
       \sum_{m\ge0}\frac{(-1)^m}{m!}\gamma_m(a)u^m.
\]
For $m,r\ge0$, put
\begin{align}
 f_{m,r}(x)&=\gamma_m^{(r)}(\{x\}),\qquad
 h_{m,r}(z)=\gamma_m^{(r)}(1+z),\label{gcx:fg}\\
 d_r&=r+1,\qquad
 L_{m,r}(X)=\prod_{j=1}^{r}(\partial_X-j)X^m.
 \label{gcx:sing-polynomial}
\end{align}
The empty product is the identity.  Thus $L_{m,0}(X)=X^m$ and
$L_{0,r}(X)=(-1)^r r!$.  The exact one-sided local decomposition is
\begin{equation}\label{gcx:local}
 f_{m,r}(x)=\boldsymbol 1_{x>0}\,
              x^{-r-1}L_{m,r}(\log x)+h_{m,r}(x),
 \qquad -1<x<1,\quad x\ne0.
\end{equation}
Derivatives in this formula are ordinary derivatives on the open
intervals.  They are not derivatives of already extended periodic
distributions.

For completeness, the singular polynomial has the harmonic form
\begin{equation}\label{gcx:harmonic-poly}
 L_{m,r}(X)=(-1)^r r!
 \sum_{j=0}^{\min(m,r)}(-1)^j\binom mj j!
 e_j\left(1,\frac12,\ldots,\frac1r\right)X^{m-j},
\end{equation}
where $e_j$ is an elementary symmetric polynomial and $e_0=1$.
Indeed,
\[
 L_{m,r}(X)=(-1)^r m![z^m](1-z)_r e^{zX}.
\]
All analytic tails in \eqref{gcx:local} are holomorphic near zero.  Their
Taylor coefficients require only ordinary Stieltjes constants and zeta
jets: for $k\ge1$,
\begin{equation}\label{gcx:tail-derivative}
 \gamma_m^{(k)}(1)=(-1)^{m+k}m![u^m](1+u)_k\zeta(k+1+u).
\end{equation}
This follows by differentiating the Hurwitz identity
$\partial_a^k\zeta(s,a)=(-1)^k(s)_k\zeta(s+k,a)$.

Fix a periodic smooth weight $w$ which is real analytic in a neighborhood
of $0$; $w=1$ and $w(x)=e^{2\pi i n x}$ are allowed.  Fix indices
$(m_i,r_i)$, $1\le i\le d$.  For distinct sufficiently small real shifts
$s_i$, define
\begin{equation}\label{gcx:CQ}
 C_{\boldsymbol m,\boldsymbol r}^{w}(\boldsymbol s)
 =\FP\int_{\R/\Z} w(x)\prod_{i=1}^d f_{m_i,r_i}(x+s_i)\,\dd x,
 \qquad
 Q_{\boldsymbol m,\boldsymbol r}^{w}
 =\FP\int_0^1w(x)\prod_{i=1}^d\gamma_{m_i}^{(r_i)}(x)\,\dd x.
\end{equation}
At each singular point $a$, the finite part removes $(a,a+\eta)$ and
retains the constant term in the ambient coordinate $\eta$.  Distinct
points use independent cutoffs.  The same convention is used at the
merged endpoint.  In particular, no cutoff is divided by a shrinking
separation.

\subsection{A finite Hermite construction}

For a nonempty subset $S\subseteq\{1,\ldots,d\}$, write
\begin{align}
 d_i&=r_i+1,\quad N_S=\sum_{i\in S}d_i,\quad
 b_S=\min_{i\in S}s_i,\notag\\
 q_S(x;\boldsymbol s)&=w(x)\prod_{j\notin S}h_{m_j,r_j}(x+s_j),
 \qquad
 D_S(x;\boldsymbol s)=\prod_{i\in S}(x+s_i)^{d_i},
 \label{gcx:qD}\\
 \Lambda_S(x;\boldsymbol s)&=
              \prod_{i\in S}L_{m_i,r_i}(\log(x+s_i)),
 \qquad x>-b_S.
 \label{gcx:Lambda}
\end{align}
The polynomial $D_S(x;\boldsymbol s)$ is monic of degree $N_S$.  Let $P_S(x;\boldsymbol s)$ be the unique polynomial of degree
less than $N_S$ that agrees with $q_S$ at $x=-s_i$ through derivative
order $d_i-1$, for every $i\in S$.

\begin{definition}[Canonical logarithmic collision counterterm]
\label{gcx:counterterm-definition}
For $|S|\ge2$, define
\begin{equation}\label{gcx:T}
 T_S(\boldsymbol s)=\FP_{R\to\infty}\FP_{\eta\downarrow0}
 \int_{-b_S+\eta}^{R}
       \frac{P_S(x;\boldsymbol s)}{D_S(x;\boldsymbol s)}
                  \Lambda_S(x;\boldsymbol s)\,\dd x,
 \qquad
 T(\boldsymbol s)=\sum_{|S|\ge2}T_S(\boldsymbol s).
\end{equation}
The lower finite part is taken first, with the ambient displacement
$\eta$.  At infinity, discard the nonconstant polynomial in $\log R$
and retain its constant term.  These are two separate endpoint
operations, not a diagonal spectral regularization.
\end{definition}

The upper operation in \eqref{gcx:T} is explicit.  Put
\[
 p_S(\boldsymbol s)=[x^{N_S-1}]P_S(x;\boldsymbol s),\qquad
 \lambda_S(X)=\prod_{i\in S}L_{m_i,r_i}(X),\qquad
 \Psi_S(X)=\int_0^X\lambda_S(t)\,\dd t.
\]
The only divergent upper term is
$p_S(\boldsymbol s)\Psi_S(\log R)$.  The difference between its derivative
and the integrand is $O(R^{-2}(1+\log R)^{\sum_{i\in S}m_i})$.
For a singleton $S$, the same definition gives zero, since translation
to its sole endpoint leaves a sum of pure monomials
$x^{-k}(\log x)^j$, $k\ge1$, whose two-endpoint finite parts vanish.

\begin{theorem}[All indices, all ordinary derivatives, and analytic weights]
\label{gcx:main}
For the convention \eqref{gcx:CQ}, the counterterm \eqref{gcx:T} satisfies
\begin{equation}\label{gcx:limit}
 C_{\boldsymbol m,\boldsymbol r}^{w}(\boldsymbol s)
       -T(\boldsymbol s)
       \longrightarrow Q_{\boldsymbol m,\boldsymbol r}^{w}
 \qquad (\boldsymbol s\longrightarrow0)
\end{equation}
through distinct real shifts.  The convergence imposes no lower bound on
ratios of the gaps.  In particular, it holds for arbitrary hierarchical
collisions and for paths whose relative gaps vanish faster than every
power of their overall size.

Every $T_S$ is a finite combination of rational functions of the gaps,
values and derivatives of $w$ and the analytic Stieltjes tails at the
shifted nodes, and generalized polylogarithms of weight at most
$1+\sum_{i\in S}m_i$.  All logarithms in its defining real integral have
positive arguments.  If every $m_i=0$, the formula is exactly the
previous rational--logarithmic collision counterterm, including its
ordinary-derivative extension.
\end{theorem}

\begin{proof}[Proof of the collision limit]
Choose $\delta>0$ so small that $w$ and all required tails are holomorphic
on $|z|<3\delta$, and restrict $|s_i|<\delta/8$.  Use the circle coordinate
$(-\delta,1-\delta)$ and separate the local interval
$(-\delta,\delta)$.  The outside integral is continuous in the shifts
(and analytic when the outside weight is analytic).  The local product
expands by the subset $S$ of factors contributing their singular terms in
\eqref{gcx:local}.  The empty subset is analytic.  The nonempty subset is
\begin{equation}\label{gcx:local-subset}
 \FP\int_{-b_S}^{\delta}
        \frac{q_S(x;\boldsymbol s)}{D_S(x;\boldsymbol s)}
                  \Lambda_S(x;\boldsymbol s)\,\dd x.
\end{equation}
Only the endpoint $-b_S$ is singular in this integral: all other nodes
$-s_i$ of the subset lie to its left.

Both $P_S$ and the divided remainder
\[
 A_S(x;\boldsymbol s)=\frac{q_S(x;\boldsymbol s)-P_S(x;\boldsymbol s)}
                           {D_S(x;\boldsymbol s)}
\]
extend holomorphically through coalescing nodes.  This assertion is not
based on individually singular Lagrange coefficients.  It follows from
the fixed-contour formulas
\begin{align}
 P_S(x;\boldsymbol s)&=\frac1{2\pi i}\oint_{|z|=2\delta}
 \frac{q_S(z;\boldsymbol s)}{D_S(z;\boldsymbol s)}
 \frac{D_S(z;\boldsymbol s)-D_S(x;\boldsymbol s)}{z-x}\,\dd z,
 \label{gcx:Hermite-contour}\\
 A_S(x;\boldsymbol s)&=\frac1{2\pi i}\oint_{|z|=2\delta}
 \frac{q_S(z;\boldsymbol s)}{D_S(z;\boldsymbol s)(z-x)}\,\dd z,
 \qquad |x|<\delta.
 \label{gcx:remainder-contour}
\end{align}
These also give a uniform bound for $A_S$ on the local integration
interval after decreasing $\delta$ or the shift neighborhood slightly.

The finite-part tail of the principal part has the convergent form
\begin{align}
 B_S(\boldsymbol s):={}&\FP\int_\delta^\infty
      \frac{P_S(x;\boldsymbol s)}{D_S(x;\boldsymbol s)}
                  \Lambda_S(x;\boldsymbol s)\,\dd x\notag\\
 ={}&-p_S(\boldsymbol s)\Psi_S(\log\delta)
 +\int_\delta^\infty\left\{
      \frac{P_S(x;\boldsymbol s)}{D_S(x;\boldsymbol s)}
                    \Lambda_S(x;\boldsymbol s)
       -\frac{p_S(\boldsymbol s)}x\lambda_S(\log x)\right\}\,\dd x.
 \label{gcx:tail}
\end{align}
For $x\ge\delta$ use the holomorphic branches of $\Log(x+s_i)$ which are
real at real shifts.  The integrand and its fixed parameter derivatives
are locally uniformly $O(x^{-2}(1+\log x)^M)$, with
$M=\sum_{i\in S}m_i$.  Consequently $B_S$ is holomorphic near the full
diagonal.

Subtracting $T_S$ from \eqref{gcx:local-subset} leaves exactly
\begin{equation}\label{gcx:after-subtraction}
 \int_{-b_S}^{\delta}A_S(x;\boldsymbol s)
                \Lambda_S(x;\boldsymbol s)\,\dd x-B_S(\boldsymbol s).
\end{equation}
The rational denominators have disappeared from the remaining local
integral.  To prove the required uniformity, set $y=x+b_S$ and
$a_i=s_i-b_S\ge0$.  Its integrand becomes
\[
 A_S(y-b_S;\boldsymbol s)
       \prod_{i\in S}L_{m_i,r_i}(\log(y+a_i)).
\]
For $\delta$ small, all $y+a_i$ in this integral are less than $1$.
Because $y+a_i\ge y$, its absolute value is bounded by
$K(1+|\log y|)^M$, independently of every gap ratio.  This is integrable
at $y=0$.  On $y\ge\epsilon>0$ the convergence is uniform as all shifts
vanish; the integral over $0<y<\epsilon$ is uniformly small as
$\epsilon\downarrow0$.  Hence \eqref{gcx:after-subtraction} tends to
\[
 \int_0^\delta A_S(x;0)\lambda_S(\log x)\,\dd x-B_S(0).
\]
Here $P_S(x;0)$ is the Taylor polynomial of $q_S(x;0)$ through degree
$N_S-1$.  The two-endpoint finite part of each monomial
$x^{j-N_S}(\log x)^k$, $j<N_S$, is zero, so
\[
 -B_S(0)=\FP\int_0^\delta
          P_S(x;0)x^{-N_S}\lambda_S(\log x)\,\dd x.
\]
Thus the limit is precisely the directly merged finite part of the
subset contribution.  Summing the finitely many subsets and the outside
integral proves \eqref{gcx:limit}.  The proof did not fix an ordering
chamber; there are only finitely many choices for the minimum, and the
same dominating function applies to all of them.
\end{proof}

\begin{proof}[Proof of finite polylogarithmic reduction]
For distinct nodes, the rational principal part is
\begin{equation}\label{gcx:partial-fractions}
 \frac{P_S(x;\boldsymbol s)}{D_S(x;\boldsymbol s)}
  =\sum_{i\in S}\sum_{k=1}^{d_i}
       \frac{b_{S,i,k}(\boldsymbol s)}{(x+s_i)^k},
\end{equation}
where
\begin{equation}\label{gcx:b-coefficient}
 b_{S,i,k}=\frac1{(d_i-k)!}
 \left.\partial_x^{d_i-k}
 \frac{q_S(x;\boldsymbol s)}
      {\prod_{j\in S\setminus\{i\}}(x+s_j)^{d_j}}
 \right|_{x=-s_i}.
\end{equation}
These are finite expressions in the stated tail and weight data.
Translate to $y=x+b_S$.  All letters $-a_i$, where
$a_i=s_i-b_S\ge0$, lie on the nonpositive real axis.  For $a>0$,
\[
 \log(y+a)=\log a+G(-a;y),\qquad
 G(-a;y)=\int_0^y\frac{\dd t}{t+a},
\]
while $G(0;y)=\log y$ uses its usual tangential normalization.
The standard iterated-integral shuffle rule \cite{goncharov} expands each logarithmic product into words of
length at most $M$.  A simple pole in \eqref{gcx:partial-fractions}
prepends one letter and therefore gives weight at most $M+1$.

For a pole of order $k\ge2$, integration by parts gives a boundary term
of logarithmic degree at most $M$ and rational integrals of logarithmic
degree at most $M-1$.  Repeat after partial fractions; induction on $M$
terminates, even if a repeated pole is transferred to another letter.
This proves the same weight bound.  Lower and upper finite parts are the
constant terms of the resulting iterated integrals at $0$ and $\infty$;
the real integral in \eqref{gcx:T} fixes their branches without crossing
any interior letter.

When all $m_i=0$, every logarithmic polynomial is constant.  The upper
primitive has zero constant term, and evaluation at the lower endpoint
gives precisely the sums of logarithms and reciprocal powers in the
index-zero repeated-pole formula.  This proves exact compatibility with
the previous convention.
\end{proof}

\begin{remark}[What regularity has been proved]\label{gcx:regularity}
The result asserts continuity after the full subtraction at arbitrary
Stieltjes indices.  The residual logarithmic integral need not be
analytic in the separations.  When all $m_i=0$, its integrand is analytic
and the lower endpoint is an analytic coordinate on each ordering
chamber; then the earlier stronger analytic-extension statement is
recovered.  No differentiability of the higher-index remainder is used
to manufacture contact terms.
\end{remark}
